\scanpage{21}

Considérons sur $\mathcal A_Q$ la forme quadratique $N_Q$ définie par \deleted{\unsure{la}}

\[
N_Q(\deleted{\unsure{u}}x)=N(u'_Q(x))
\]

on aura donc, en désignant par $1_Q$ la section $1$ de $\underline O_{\unsure{\alpha|X}}$, regardée comme section de $\mathcal A_Q$

\[
N_Q(1_Q)=1,
\]

Ainsi,

\[
\deleted{\unsure{A}}(\mathcal A_Q,1_Q,N_Q)
\]

est un module loc. libre de rang 2, muni d'une section $1_Q$ et d'une forme quadratique valant $1$ sur cette section. On a donc une structure d'alg. quadratique unique sur $\mathcal A_Q$, pour laquelle $N_Q$ soit la norme. De plus, on vérifie immédiatement que $u'_Q$ est un hom d'algèbres, et qu'on a, en plus de la relation sur les normes, la relation

\[
\begin{aligned}
\deleted{\unsure{N}}\operatorname{Tr}(u'_Q(x))&=\operatorname{Tr}_{\mathcal A_Q/\underline O}(x)\\
u'_Q(\bar x)&=u'_Q(x)^{\text{--}}
\end{aligned}
\]

\unsure{ici|où} on désigne par

\[
v\longmapsto\bar v=(\operatorname{Tr}v)\cdot1-v
\]

l'antiinvolution canonique de $\underline{\operatorname{End}}(E)$, égale à $\mathrm{id}_{\underline O}$ sur $\underline O\cdot1_E$ et induisant $-\mathrm{id}$ sur le quotient.
% archive sheet number: 20
